\section{AHSME 1991}
        \begin{enumerate}
        	\item (\textit{AHSME 1991}) If for any three distinct numbers $a$, $b$, and $c$ we define $f(a,b,c)=\frac{c+a}{c-b}$, then $f(1,-2,-3)$ is


 $\textbf{(A) } -2 \qquad \textbf{(B) } -\frac{2}{5} \qquad \textbf{(C) } -\frac{1}{4} \qquad \textbf{(D) } \frac{2}{5} \qquad \textbf {(E) } 2$


 

	\item (\textit{AHSME 1991}) $|3-\pi|=$


 $\textbf{(A) }\frac{1}{7} \qquad \textbf{(B) }0.14 \qquad \textbf{(C) }3-\pi \qquad \textbf{(D) }3+\pi \qquad \textbf{(E) }\pi-3$


 

	\item (\textit{AHSME 1991}) $(4^{-1}-3^{-1})^{-1}=$


 $\textbf{(A) }-12 \qquad \textbf{(B) }-1 \qquad \textbf{(C) }\frac{1}{12} \qquad \textbf{(D) }1 \qquad \textbf{(E) }12$


 

	\item (\textit{AHSME 1991}) Which of the following triangles cannot exist?


 $\textbf{(A) }$ An acute isosceles triangle


 $\textbf{(B) }$ An isosceles right triangle


 $\textbf{(C) }$ An obtuse right triangle


 $\textbf{(D) }$ A scalene right triangle


 $\textbf{(E) }$ A scalene obtuse triangle


 

	\item (\textit{AHSME 1991}) 
                    \begin{figure}[H]
                        \centering
                        \includegraphics[width=0.5\textwidth]{Problem/AHSME_I_1991_5_Problem_Figure_1.png}
                    \end{figure}
                    


 In the arrow-shaped polygon [see figure], the angles at vertices $A,C,D,E$ and $F$ are right angles, $BC=FG=5, CD=FE=20, DE=10$, and $AB=AG$. The area of the polygon is closest to


 $\textbf{(A) } 288\qquad\textbf{(B) } 291\qquad\textbf{(C) } 294\qquad\textbf{(D) } 297\qquad\textbf{(E) } 300$


 

	\item (\textit{AHSME 1991}) If $x\geq 0$, then $\sqrt{x\sqrt{x\sqrt{x}}}=$


 $\textbf{(A) } x\sqrt{x}\qquad \textbf{(B) } x\sqrt[4]{x}\qquad \textbf{(C) } \sqrt[8]{x}\qquad \textbf{(D) } \sqrt[8]{x^3}\qquad \textbf{(E) } \sqrt[8]{x^7}$


 

	\item (\textit{AHSME 1991}) If $x=\frac{a}{b}$, $a\neq b$ and $b\neq 0$, then $\frac{a+b}{a-b}=$


 $\textbf{(A) } \frac{x}{x+1} \qquad \textbf{(B) } \frac{x+1}{x-1} \qquad \textbf{(C) } 1 \qquad \textbf{(D) } x-\frac{1}{x} \qquad \textbf{(E) } x+\frac{1}{x}$


 

	\item (\textit{AHSME 1991}) Liquid $X$ does not mix with water. Unless obstructed, it spreads out on the surface of water to form a circular film $0.1$cm thick. A rectangular box measuring $6$cm by $3$cm by $12$cm is filled with liquid $X$. Its contents are poured onto a large body of water. What will be the radius, in centimeters, of the resulting circular film?


 $\textbf{(A) } \frac{\sqrt{216}}{\pi} \qquad \textbf{(B) }\sqrt{\frac{216}{\pi}} \qquad \textbf{(C) } \sqrt{\frac{2160}{\pi}} \qquad \textbf{(D) } \frac{216}{\pi} \qquad \textbf{(E) } \frac{2160}{\pi}$


 

	\item (\textit{AHSME 1991}) From time $t=0$ to time $t=1$ a population increased by $i\%$, and from time $t=1$ to time $t=2$ the population increased by $j\%$. Therefore, from time $t=0$ to time $t=2$ the population increased by


 $\textbf{(A) } (i+j)\% \qquad \textbf{(B) } ij\% \qquad \textbf{(C) } (i+ij)\% \qquad \textbf{(D) } \left(i+j+\frac{ij}{100}\right)\% \qquad \textbf{(E) } \left(i+j+\frac{i+j}{100}\right)\%$


 

	\item (\textit{AHSME 1991}) Point $P$ is $9$ units from the center of a circle of radius $15$. How many different chords of the circle contain $P$ and have integer lengths?


 $\textbf{(A) } 11\qquad \textbf{(B) } 12\qquad \textbf{(C) } 13\qquad \textbf{(D) } 14\qquad \textbf{(E) } 29$


 

	\item (\textit{AHSME 1991}) Jack and Jill run 10 km. They start at the same point, run 5 km up a hill, and return to the starting point by the same route. Jack has a 10 minute head start and runs at the rate of 15 km/hr uphill and 20 km/hr downhill. Jill runs 16 km/hr uphill and 22 km/hr downhill. How far from the top of the hill are they when they pass each other going in opposite directions (in km)?


 $\textbf{(A) } \frac{5}{4}\qquad \textbf{(B) } \frac{35}{27}\qquad \textbf{(C) } \frac{27}{20}\qquad \textbf{(D) } \frac{7}{3}\qquad \textbf{(E) } \frac{28}{49}$


 

	\item (\textit{AHSME 1991}) The measures (in degrees) of the interior angles of a convex hexagon form an arithmetic sequence of integers. Let $m$ be the measure of the largest interior angle of the hexagon. The largest possible value of $m$, in degrees, is


 $\textbf{(A) } 165\qquad \textbf{(B) } 167\qquad \textbf{(C) } 170\qquad \textbf{(D) } 175\qquad \textbf{(E) } 179$


 

	\item (\textit{AHSME 1991}) Horses $X,Y$ and $Z$ are entered in a three-horse race in which ties are not possible. The odds against $X$ winning are $3:1$ and the odds against $Y$ winning are $2:3$, what are the odds against $Z$ winning? (By "odds against $H$ winning are $p:q$" we mean the probability of $H$ winning the race is $\frac{q}{p+q}$.)


 $\textbf{(A) } 3:20\qquad \textbf{(B) } 5:6\qquad \textbf{(C) } 8:5\qquad \textbf{(D) } 17:3\qquad \textbf{(E) } 20:3$


 

	\item (\textit{AHSME 1991}) If $x$ is the cube of a positive integer and $d$ is the number of positive integers that are divisors of $x$, then $d$ could be


 $\textbf{(A) } 200\qquad \textbf{(B) } 201\qquad \textbf{(C) } 202\qquad \textbf{(D) } 203\qquad \textbf{(E) } 204$


 

	\item (\textit{AHSME 1991}) A circular table has 60 chairs around it. There are $N$ people seated at this table in such a way that the next person seated must sit next to someone. What is the smallest possible value for $N$?


 $\textbf{(A) } 15\qquad \textbf{(B) } 20\qquad \textbf{(C) } 30\qquad \textbf{(D) } 40\qquad \textbf{(E) } 58$


 

	\item (\textit{AHSME 1991}) One hundred students at Century High School participated in the AHSME last year, and their mean score was 100. The number of non-seniors taking the AHSME was $50\%$ more than the number of seniors, and the mean score of the seniors was $50\%$ higher than that of the non-seniors. What was the mean score of the seniors?


 $\textbf{(A) } 100\qquad \textbf{(B) } 112.5\qquad \textbf{(C) } 120\qquad \textbf{(D) } 125\qquad \textbf{(E) } 150$


 

	\item (\textit{AHSME 1991}) A positive integer $N$ is a \textit{palindrome} if the integer obtained by reversing the sequence of digits of $N$ is equal to $N$. The year 1991 is the only year in the current century with the following 2 properties:


 (a) It is a palindrome
(b) It factors as a product of a 2-digit prime palindrome and a 3-digit prime palindrome. 


 How many years in the millenium between 1000 and 2000 have properties (a) and (b)?


 $\textbf{(A) } 1\qquad \textbf{(B) } 2\qquad \textbf{(C) } 3\qquad \textbf{(D) } 4\qquad \textbf{(E) } 5$


 

	\item (\textit{AHSME 1991}) If $S$ is the set of points $z$ in the complex plane such that $(3+4i)z$ is a real number, then $S$ is a 


 $\textbf{(A) }$ right triangle $\qquad \textbf{(B) }$ circle $\qquad \textbf{(C) }$ hyperbola $\qquad \textbf{(D) }$ line $\qquad \textbf{(E) }$ parabola


 

	\item (\textit{AHSME 1991}) 
                    \begin{figure}[H]
                        \centering
                        \includegraphics[width=0.5\textwidth]{Problem/AHSME_I_1991_19_Problem_Figure_1.png}
                    \end{figure}
                    


 Triangle $ABC$ has a right angle at $C, AC=3$ and $BC=4$. Triangle $ABD$ has a right angle at $A$ and $AD=12$. Points $C$ and $D$ are on opposite sides of $\overline{AB}$. The line through $D$ parallel to $\overline{AC}$ meets $\overline{CB}$ extended at $E$. If

 \[\frac{DE}{DB}=\frac{m}{n},\]
where $m$ and $n$ are relatively prime positive integers, then $m+n$ is


 $\textbf{(A) } 25\qquad \textbf{(B) } 128\qquad \textbf{(C) } 153\qquad \textbf{(D) } 243\qquad \textbf{(E) } 256$


 

	\item (\textit{AHSME 1991}) The sum of all real $x$ such that $(2^x-4)^3+(4^x-2)^3=(4^x+2^x-6)^3$ is


 $\textbf{(A) } \frac32 \qquad \textbf{(B) } 2 \qquad \textbf{(C) } \frac52 \qquad \textbf{(D) } 3 \qquad \textbf{(E) } \frac72$


 

	\item (\textit{AHSME 1991}) For all real numbers $x$ except $x=0$ and $x=1$ the function $f(x)$ is defined by $f(x/(x-1))=1/x$. Suppose $0\leq t\leq \pi/2$. What is the value of $f(\sec^2t)$?


 $\textbf{(A) } \sin^2\theta\qquad \textbf{(B) } \cos^2\theta\qquad \textbf{(C) } \tan^2\theta\qquad \textbf{(D) } \cot^2\theta\qquad \textbf{(E) } \csc^2\theta$


 

	\item (\textit{AHSME 1991}) 
                    \begin{figure}[H]
                        \centering
                        \includegraphics[width=0.5\textwidth]{Problem/AHSME_I_1991_22_Problem_Figure_1.png}
                    \end{figure}
                    


 Two circles are externally tangent. Lines $\overline{PAB}$ and $\overline{PA'B'}$ are common tangents with $A$ and $A'$ on the smaller circle $B$ and $B'$ on the larger circle. If $PA=AB=4$, then the area of the smaller circle is


 $\textbf{(A) } 1.44\pi\qquad \textbf{(B) } 2\pi\qquad \textbf{(C) } 2.56\pi\qquad \textbf{(D) } \sqrt{8}\pi\qquad \textbf{(E) } 4\pi$


 

	\item (\textit{AHSME 1991}) 
                    \begin{figure}[H]
                        \centering
                        \includegraphics[width=0.5\textwidth]{Problem/AHSME_I_1991_23_Problem_Figure_1.png}
                    \end{figure}
                    


 If $ABCD$ is a $2\times2$ square, $E$ is the midpoint of $\overline{AB}$,$F$ is the midpoint of $\overline{BC}$,$\overline{AF}$ and $\overline{DE}$ intersect at $I$, and $\overline{BD}$ and $\overline{AF}$ intersect at $H$, then the area of quadrilateral $BEIH$ is


 $\textbf{(A) } \frac{1}{3}\qquad \textbf{(B) } \frac{2}{5}\qquad \textbf{(C) } \frac{7}{15}\qquad \textbf{(D) } \frac{8}{15}\qquad \textbf{(E) } \frac{3}{5}$


 

	\item (\textit{AHSME 1991}) The graph, $G$ of $y=\log_{10}x$ is rotated $90^{\circ}$ counter-clockwise about the origin to obtain a new graph $G'$. Which of the following is an equation for $G'$?


 $\textbf{(A) } y=\log_{10}\left(\frac{x+90}{9}\right) \qquad \textbf{(B) } y=\log_{x}10 \qquad \textbf{(C) } y=\frac{1}{x+1} \qquad \textbf{(D) } y=10^{-x} \qquad \textbf{(E) } y=10^x$


 

	\item (\textit{AHSME 1991}) If $T_n=1+2+3+\cdots +n$ and

 \[P_n=\frac{T_2}{T_2-1}\cdot\frac{T_3}{T_3-1}\cdot\frac{T_4}{T_4-1}\cdot\cdots\cdot\frac{T_n}{T_n-1}\]
for $n=2,3,4,\cdots,$ then $P_{1991}$ is closest to which of the following numbers?


 $\textbf{(A) } 2.0\qquad \textbf{(B) } 2.3\qquad \textbf{(C) } 2.6\qquad \textbf{(D) } 2.9\qquad \textbf{(E) } 3.2$


 

	\item (\textit{AHSME 1991}) An $n$-digit positive integer is cute if its $n$ digits are an arrangement of the set $\{1,2,...,n\}$ and its first 
$k$ digits form an integer that is divisible by $k$  , for  $k  = 1,2,...,n$. For example, $321$ is a cute $3$-digit integer because $1$ divides $3$, $2$ divides $32$, and $3$ divides $321$. How many cute $6$-digit integers are there?


 $\textbf{(A) } 0\qquad \textbf{(B) } 1\qquad \textbf{(C) } 2\qquad \textbf{(D) } 3\qquad \textbf{(E) } 4$


 

	\item (\textit{AHSME 1991}) If 
 \[x+\sqrt{x^2-1}+\frac{1}{x-\sqrt{x^2-1}}=20,\] then 
 \[x^2+\sqrt{x^4-1}+\frac{1}{x^2+\sqrt{x^4-1}}=\]


 $\textbf{(A) } 5.05 \qquad \textbf{(B) } 20 \qquad \textbf{(C) } 51.005 \qquad \textbf{(D) } 61.25 \qquad \textbf{(E) } 400$


 

	\item (\textit{AHSME 1991}) Initially an urn contains 100 white and 100 black marbles. Repeatedly 3 marbles are removed (at random) from the urn and replaced with some marbles from a pile outside the urn as follows: 3 blacks are replaced with 1 black, or 2 blacks and 1 white are replaced with a white and a black, or 1 black and 2 whites are replaced with 2 whites, or 3 whites are replaced with a black and a white. Which of the following could be the contents of the urn after repeated applications of this procedure?


 $\textbf{(A) }$ 2 black $\qquad \textbf{(B) }$ 2 white $\qquad \textbf{(C) }$ 1 black $\qquad \textbf{(D) }$ 1 black and 1 white $\qquad \textbf{(E) }$ 1 white


 

	\item (\textit{AHSME 1991}) Equilateral triangle $ABC$ has $P$ on $AB$ and $Q$ on $AC$. The triangle is folded along $PQ$ so that vertex $A$ now rests at $A'$ on side $BC$. If $BA'=1$ and $A'C=2$ then the length of the crease $PQ$ is


 $\textbf{(A) } \frac{8}{5} \qquad \textbf{(B) } \frac{7}{20}\sqrt{21} \qquad \textbf{(C) } \frac{1+\sqrt{5}}{2} \qquad \textbf{(D) } \frac{13}{8} \qquad \textbf{(E) } \sqrt{3}$


 

	\item (\textit{AHSME 1991}) For any set $S$, let $|S|$ denote the number of elements in $S$, and let $n(S)$ be the number of subsets of $S$, including the empty set and the set $S$ itself. If $A$, $B$, and $C$ are sets for which $n(A)+n(B)+n(C)=n(A\cup B\cup C)$ and $|A|=|B|=100$, then what is the minimum possible value of $|A\cap B\cap C|$?


 $\textbf{(A) } 96 \qquad \textbf{(B) } 97 \qquad \textbf{(C) } 98 \qquad \textbf{(D) } 99 \qquad \textbf{(E) } 100$


 

\end{enumerate}